\clearpage\begingroup\SourceLeaf{77b}\NoSourceHead
\begin{center}[2.]\end{center}
\begin{center}\textsc{Gauss} an \textsc{Olbers}.\quad Göttingen, 18. Februar 1832.\end{center}
\begin{center}
\rule{.70\textwidth}{.4pt}\\
{\small [\textsc{Wilhelm Olbers}, Sein Leben und seine Werke II, 2, 1909, S. 584.]}\\[-.3ex]
\rule{.70\textwidth}{.4pt}
\end{center}
\noindent . . . . . . . . . . . . . . . . . . . . .\par
\indent Ich beschäftige mich jetzt mit dem Erdmagnetismus, namentlich mit einer
\textls[120]{absoluten} Bestimmung von dessen Intensität. Freund \textsc{Weber} macht nach
meiner Angabe die Versuche. So wie man z. B. von Geschwindigkeit nur
durch Ansetzung einer Zeit und eines Raums einen klaren Begriff geben kann,
so, finde ich, muss zur vollständigen Bestimmung der Intensität des Erdmagnetismus
angegeben werden 1) ein Gewicht \(=p\), 2) eine Linie \(=r\), und dann
kann man den Erdmagnetismus durch \(\dfrac{\sqrt p}{r}\) ausdrücken, d. i. ein doppelt ge-
\clearpage\SourceLeaf{78a}\HistoricalHead{73}{R}{GAUSS AN OLBERS.}
\noindent wordener Erdmagnetismus, würde bei gleichem \(r\) ein viermahl so grosses Gewicht,
oder bei gleichem Gewicht ein halbsogrosses \(r\) herbeiführen. Es
scheint, dass, wenn man für \(r\) einen Zoll nimmt, \(p\) wohl nur wenige Milligramm
beträgt. Die Versuche sind aber noch nicht vollständig. Ich werde,
wenn es Sie interessirt, Ihnen gern demnächst etwas Näheres mittheilen und
bemerke nur, dass man dabei \textls[120]{zwei} Nadeln \(A\) und \(B\) nöthig hat (die eine ist
übrigens ein Stab), dass die Wirkung des Erdmagnetismus auf \(B\) mit der
Wirkung von \(A\) auf \(B\) vergleichbar ist, insofern man letztere in bestimmter
nicht zu kleiner Entfernung spielen lässt, deren Cubus die letztere Wirkung
umgekehrt proportional ist; die Wirkung des Erdmagnetismus auf \(A\) hingegen
ist mit dem Momente eines Gewichts, Produkt des Letzteren in eine Linie,
vergleichbar, was dann entweder durch die Wage, indem man ein kleines
Gewicht jene Wirkung aufheben lässt, oder durch Beobachtung der Schwingungszeiten
ausgemittelt werden kann.

Auch für Declination und Inclination hoffe ich, mehrere neue Verbesserungen
mit \textsc{Webers} Hilfe angeben zu können. . . . . . . . . . . . . . . .
\clearpage\endgroup
